%% =====================================================================
%%  B8-T-23-09 -- what B8 contributes to "The Fifty-Four Percent Ceiling"
%%  -------------------------------------------------------------------
%%  LAYER A: APPEND-ONLY. Written once at the entry's write, then left
%%  alone. If the source has more to say later, add a bullet -- do not
%%  rewrite. Material found ONLY here goes in the `unique` slot with
%%  the unique flag set, which makes it undeletable (rule R10).
%% =====================================================================
%% @kind: source
%% @id: B8-T-23-09
%% @book: B8
%% @topic: T-23-09
%% @locator: PSPR 10(3):214--234; abstract, Results (percentage correct, determinants of accuracy), Discussion
%% @status: distilled
%% @unique: yes
%% @integrated: yes
%% @updated: 2026-09-20

\dossier{B8}{T-23-09}{\bookshort{B8}}{PSPR 10(3):214--234; abstract, Results, Discussion}

\begin{distilled}
  \item The ceiling, measured: pooling 206 documents, 24,483 judges, 4,435 senders and 6,651 messages, mean unaided real-time lie-truth discrimination is 53.98\% raw and 53.46\% weighted (95\% CI 53.31--53.59). The spread of study means is tiny (quartiles 50.07--58.00; true between-study SD 4.52 points) --- the ceiling is stable, not an artifact of weak studies.
  \item The errors are asymmetric: 61.34\% of truths are correctly accepted against 47.55\% of lies correctly rejected; the weighted truth bias is 55.23\% of all judgments calling a message true. Signal-detection measures agree (mean $d' = .24$; maximum achievable 54.79\%).
  \item The edge is real but thin: accuracy $d \approx .39$--$.40$, ranking at the 60th percentile of 474 meta-analysed social-psychological effects.
  \item The eye is the weak channel: video-only judgment discriminates at near chance (between-study $d = .077$) against audio ($d = .419$) and audiovisual ($d = .438$); face-only $d = .01$, body-only $d = -.15$.
  \item Motivation backfires into the stereotype: senders motivated to be believed were called truthful by 46.84\% of video-only receivers versus 54.44\% for unmotivated senders --- whether or not they were lying; motivation produced no within-study accuracy gain.
  \item Preparation matters at the sender's end: unprepared senders are easier to read (within-study $d = -.144$); planned messages merely appear more truthful ($d = .133$).
  \item Baseline exposure is the one receiver-side lever: judges with prior exposure to the sender reach 55.91\% against 52.26\% without it (within-study $d = .239$).
  \item The expertise null: across 19 studies where professionals and novices judged the same messages, weighted $d = -.025$ (95\% CI $-.105$ to $.055$); experts were slightly more skeptical (52.28\% vs 55.69\% truth classifications) without being more right.
  \item Interaction-partner truth bias: people judge those they interact with as more truthful than third parties judge them ($d = .26$).
  \item Real-world discovery runs late: people who uncovered lies report it took days to months and turned on physical evidence and third parties, not on behaviour displayed at the moment (B8's synthesis of Park et al. 2002).
  \item The double standard: judges carry a moralistic liar stereotype (torment, fidgeting, averted gaze) while exempting their own practical lies from the category; calm, rationalized deceit fails to match the stereotype and passes.
\end{distilled}
\begin{unique}
  \item The 54\% ceiling with its measured spread, the truth--lie asymmetry (61.34/47.55), the channel hierarchy (hearing beats seeing; face and body cues separately null), the motivation backfire, and the expertise null are unique to B8 in this corpus --- no other source quantifies the detection enterprise as a whole.
\end{unique}
